% Part: incompleteness
% Chapter: theories-computability
% Section: q-is-ce

\documentclass[../../../include/open-logic-section]{subfiles}

\begin{document}

\olfileid{inc}{tcp}{qce}

\olsection{$\Th{Q}$، \printtoken{S}{c.e.}-کامل است}


\begin{thm}
$\Th{Q}$، !!{c.e.} است اما تصمیم‌پذیر نیست. در واقع، یک مجموعهٔ
!!{c.e.} کامل است.
\end{thm}

\begin{proof}
دیدن اینکه~$\Th{Q}$، !!{c.e.} است دشوار نیست، زیرا مجموعهٔ
(رمزهای) جمله‌های~$y$ است که اثباتی چون~$x$ از~$y$ در~$\Th{Q}$
وجود دارد:
\[
Q = \Setabs{y}{\lexists[x][\Prf[\Th{Q}](x,y)]}.
\]
اما می‌دانیم~$\Prf[\Th{Q}](x,y)$ محاسبه‌پذیر (و در واقع بازگشتی
اولیه) است، و هر مجموعه‌ای که بتوان آن را به صورت بالا نوشت، c.e.
است.

گفتن اینکه این مجموعه c.e. کامل است، با گفتن اینکه
$K \leq_m Q$، که در آن
$K = \Setabs{x}{!A_x(x) \downarrow}$، هم‌ارز است. پس نشان می‌دهیم
$K$ به~$\Th{Q}$ کاهش‌پذیر است. از آنجا که محمول کلینی
$T(e,x,s)$ بازگشتی اولیه است، در~$\Th{Q}$، مثلاً با~$!A_T$،
بازنمایی‌پذیر است. آنگاه برای هر~$x$ داریم
\begin{align*}
x \in K & \lif \lexists[s][T(x,x,s)] \\
& \lif \lexists[s][(\Th{Q} \Proves !A_T(\num x,\num x, \num s))]
  \\
& \lif \Th{Q} \Proves \lexists[s][!A_T(\num x, \num x, s)].
\end{align*}
برعکس، اگر
$\Th{Q} \Proves \lexists[s][!A_T(\num x, \num x, s)]$، آنگاه در
واقع برای عدد طبیعی~$n$ای، فرمول
$!A_T(\num x, \num x, \num n)$ باید صادق باشد. اگر
$T(x,x,n)$ کاذب بود، $\Th{Q}$،
$\lnot !A_T(\num x, \num x, \num n)$ را ثابت می‌کرد، زیرا
$!A_T$، $T$ را بازنمایی می‌کند.
اما در آن صورت~$\Th{Q}$ فرمولی کاذب را ثابت می‌کرد، که تناقض است.
پس~$T(x,x,n)$ باید صادق باشد و این
$!A_x(x) \downarrow$ را نتیجه می‌دهد.

خلاصه اینکه برای هر~$x$، $x$ عضو~$K$ است اگر و تنها اگر
$\Th{Q}$، $\lexists[s][!A_T(\num x,\num x,s)]$ را ثابت کند. پس
تابع~$f$ که~$x$ را به (رمزِ) جملهٔ
$\lexists[s][!A_T(\num x, \num x, s)]$ می‌برد، کاهشی از~$K$ به
$\Th{Q}$ است.
\end{proof}

\end{document}
